\documentclass[12pt,a4paper]{article}

% --- Pachete necesare ---
\usepackage[utf8]{inputenc}
\usepackage[T1]{fontenc}
\usepackage[romanian]{babel}
\usepackage{amsmath, amssymb}
\usepackage{graphicx}
\usepackage{geometry}
\usepackage{booktabs}
\usepackage{siunitx}
\usepackage{hyperref}

\geometry{margin=2.5cm}

% --- Datele referatului ---
\title{\textbf{Determinarea lungimii de undă a luminii monocromatice\\ cu ajutorul dispozitivului Young}}
\author{Nume Prenume \\ Grupa / Seria}
\date{\today}

\begin{document}

\maketitle

\section{Scopul lucrării}
\begin{itemize}
    \item Obținerea figurii de interferență a luminii cu ajutorul unui dispozitiv de tip Young.
    \item Determinarea experimentală a lungimii de undă a unei radiații luminoase monocromatice.
\end{itemize}

\section{Teoria lucrării}
Lumina este o undă electromagnetică. În dispozitivul Young, lumina provenită de la o sursă monocromatică \(S\) traversează un paravan opac cu două fante paralele și înguste \(S_1\) și \(S_2\), aflate la distanța \(d\) una de cealaltă. Undele difractate de cele două fante sunt coerente și interferă pe un ecran aflat la distanța \(D\) de planul fantelor.

Diferența de drum optic dintre undele care ajung într-un punct \(P\) de pe ecran este:
\[
\delta = \delta_2 - \delta_1 = n\lambda,
\]
unde \(n\) este ordinul de interferență. Pentru maximele de interferență, poziția pe ecran este dată de relația:
\[
x_n = n\,\frac{\lambda D}{d}.
\]

Distanța dintre două maxime (sau minime) consecutive se numește \textbf{interfranjă} și se notează cu \(i\):
\[
i = x_{n+1} - x_n = \frac{\lambda D}{d}.
\]

Din relația de mai sus se obține formula de calcul pentru lungimea de undă:
\[
\lambda = \frac{d \cdot i}{D}.
\]

În practică, dacă se notează cu \(2l\) distanța dintre fante (adică \(d = 2l\)), formula devine:
\[
\lambda = \frac{2l \cdot i}{D},
\]
unde \(i\) este interfranja măsurată experimental[reference:0].

\section{Descrierea montajului experimental}
Montajul experimental cuprinde:
\begin{itemize}
    \item o sursă de lumină monocromatică (sau sursă de lumină albă prevăzută cu filtru monocromatic);
    \item un paravan opac cu două fante paralele, subțiri și apropiate (dispozitiv Young);
    \item un ecran sau un sistem optic de observare prevăzut cu șurub micrometric pentru măsurarea pozițiilor franjelor;
    \item o riglă gradată sau un banc optic pentru măsurarea distanței \(D\) dintre planul fantelor și ecran.
\end{itemize}

Distanța dintre fante \(d\) (sau \(2l\)) este cunoscută din construcția dispozitivului. Distanța \(D\) se măsoară cu rigla sau cu ajutorul reperelor de pe bancă.

\section{Modul de lucru}
\begin{enumerate}
    \item Se aliniază sursa de lumină, dispozitivul cu fante și ecranul astfel încât să fie coliniare.
    \item Se reglează lățimea fantei sursei și poziția ecranului până când figura de interferență este clară și contrastantă.
    \item Se măsoară distanța \(D\) dintre planul fantelor și ecran.
    \item Se măsoară pozițiile a \(N\) franje luminoase consecutive (sau a unui grup de 5–10 franje) pentru a determina interfranja \(i\).
    \item Se repetă măsurătorile pentru cel puțin 5–10 poziții diferite ale ecranului sau pentru mai multe grupuri de franje.
    \item Se calculează lungimea de undă \(\lambda\) pentru fiecare determinare și se prelucrează statistic rezultatele.
\end{enumerate}

\section{Prelucrarea datelor experimentale}
\subsection*{Calculul interfranjei}
Dacă se măsoară pozițiile a \(N\) franje consecutive, interfranja se calculează cu relația:
\[
i = \frac{x_N - x_0}{N},
\]
unde \(x_0\) și \(x_N\) sunt pozițiile primei și ultimei franje din grupul măsurat.

\subsection*{Calculul lungimii de undă}
Pentru fiecare determinare se aplică formula:
\[
\lambda = \frac{d \cdot i}{D}.
\]

\subsection*{Calculul erorilor}
Se calculează valoarea medie a lungimii de undă:
\[
\bar{\lambda} = \frac{1}{n}\sum_{k=1}^{n} \lambda_k.
\]
Eroarea pătratică medie (eroarea standard) se determină cu relația:
\[
\sigma_\lambda = \sqrt{\frac{\sum_{k=1}^{n} (\lambda_k - \bar{\lambda})^2}{n(n-1)}}.
\]
Rezultatul final se prezintă sub forma:
\[
\lambda = \bar{\lambda} \pm \sigma_\lambda.
\]

\section{Tabel de date (model)}
\begin{table}[h!]
\centering
\caption{Model de tabel pentru înregistrarea datelor experimentale}
\begin{tabular}{@{}cccccc@{}}
\toprule
Nr. & \(D\) (m) & \(N\) franje & \(x_N - x_0\) (mm) & \(i\) (mm) & \(\lambda\) (nm) \\
\midrule
1 & & & & & \\
2 & & & & & \\
3 & & & & & \\
\(\vdots\) & & & & & \\
\bottomrule
\end{tabular}
\end{table}

\section*{Observații}
\begin{itemize}
    \item Distanța dintre fante \(d\) trebuie exprimată în aceleași unități de măsură ca și interfranja \(i\) și distanța \(D\).
    \item Pentru a reduce erorile, se recomandă măsurarea unui număr cât mai mare de franje și repetarea determinărilor.
    \item În cazul luminii albe, franjele sunt colorate, iar măsurătorile se fac pentru o singură radiație monocromatică (cu filtru).
\end{itemize}

\end{document}